% TEMPLATE-MILC-v3.tex - plantilla maestra UNICA de las guias MILC v3.
%
% La usan las cuatro familias de guias, cada una con su driver:
%   · Grados 6-11          -> scripts/build-guias-g11.py
%   · Bebras               -> scripts/build-guias-bebras.py
%   · Semillero            -> scripts/build-guias-semillero.py
%   · Territorio interior  -> scripts/build-guias-territorio-interior.py
%
% Antes habia tres copias casi identicas (~400 lineas iguales, 6 distintas):
% cada mejora habia que aplicarla tres veces, y en la practica no se hacia
% (el motor de recursos vivia solo en dos de ellas). Lo que varia entre
% familias esta parametrizado en 6 placeholders que cada driver llena
% (se nombran aqui SIN su sintaxis real: el builder reemplaza por texto plano,
% tambien dentro de los comentarios, y un valor multilinea rompe el preambulo):
%
%   HEADER_IZQ        encabezado izquierdo de pagina
%   HEADER_DER        encabezado derecho de pagina
%   PORTADA_ETIQUETA  rotulo sobre el numero grande (GRADO/MOMENTO/MODULO)
%   PORTADA_SERIE     linea de serie de la portada
%   PORTADA_ID        identificador de la guia en la portada
%   PORTADA_CIRCULO   el circulito de grado (°), o vacio si no aplica
%   PDF_TITULO        titulo en texto plano (sin \textbf ni comillas TeX) para
%                     los metadatos del PDF (/Title); lo produce lib_xelatex.texto_pdf
%
% Composicion (auditoria 2026-09-03): las bandas de estacion piden espacio
% (\Needspace*) y prohiben el salto justo despues (\nopagebreak), asi no quedan
% huerfanas al pie; las cajas largas (softbox, drawbox, infoband, puentebox) son
% `breakable`, asi una caja de media pagina ya no arrastra todo a la siguiente
% dejando la anterior al 40 %. Todo texto sobre mostaza o verde va en milcNegro
% (blanco sobre #E5B400 da 1,93:1 y sobre #58A923 2,95:1; ninguno pasa WCAG AA).
% El resto de claves las documenta content/guias/_SCHEMA.md. El script valida
% que no queden placeholders sin reemplazar antes de escribir el archivo final.

% Etiquetado PDF (latex-lab): /Lang, estructura y texto alternativo de las
% figuras, para lectores de pantalla. Debe ir ANTES de \documentclass.
\DocumentMetadata{lang=es-CO,tagging=on}
\documentclass[11pt,letterpaper]{article}
% headheight=24pt: la cabecera derecha lleva dos lineas (identidad de la guia
% y estacion actual). headsep se acorta en la misma medida para que el bloque
% de texto quede exactamente donde estaba con headheight=12pt.
\usepackage[margin=2.0cm,headheight=24pt,headsep=13pt]{geometry}
\usepackage[table]{xcolor}
\usepackage{fontspec}
\usepackage[spanish]{babel}
\usepackage{tikz}
% Librerías TikZ para los diagramas de las guías (bloque `recursos:`):
%  · babel  → babel en español vuelve activos caracteres como «>», y TikZ los usa
%             en sus opciones (>=stealth, ->). Sin esta librería, cualquier
%             diagrama con flechas revienta con «\language@active@arg> has an
%             extra }». Debe ir ANTES de que se dibuje nada.
%  · positioning        → sintaxis «below left=2pt and 3pt».
%  · decorations.pathreplacing → llaves ({brace}) para señalar rangos.
\usetikzlibrary{babel,positioning,decorations.pathreplacing}
\usepackage{graphicx}
% Circuitos: el trabajo de electrónica se hace hoy con micro:bit y sus sensores
% integrados (MakeCode, sin cableado), así que no se necesitan esquemáticos.
% Si algún día entran montajes con protoboard, basta descomentar esta línea:
% los diagramas .tex/.tikz declarados en `recursos:` ya se incrustan solos.
% \usepackage{circuitikz}
\usepackage[most]{tcolorbox}
\usepackage{tabularx}
\usepackage{array}
\usepackage{fancyhdr}
\usepackage{titlesec}
\usepackage[document]{ragged2e}  % bandera a la derecha en todo el cuerpo
\usepackage{enumitem}
\usepackage{microtype}
\usepackage{etoolbox}
\usepackage{needspace}
% hyperref al final del preambulo: metadatos del PDF (titulo, autor, idioma,
% familia), marcadores por estacion (\pdfbookmark) y enlaces sin recuadro.
\usepackage[hidelinks,
  pdftitle={Accesibilidad lectora — quién NO puede leer mi revista y por qué},
  pdfauthor={PhD. Álvaro Cárdenas Orozco},
  pdfsubject={Tecnología e Informática · Grado 9 · Guía 17 · MILC v3},
  pdflang={es-CO},
  bookmarksopen=true,bookmarksnumbered=false]{hyperref}

% Helvetica Neue no trae la flecha U+2192, asi que cada «→» escrito literalmente
% en el YAML se compilaba a NADA, en silencio (462 flechas en 61 guias: rutas del
% tipo «paleta → tipografia → lockup» salian rotas). Mapeamos el caracter a su
% version matematica, que si tiene glifo. Asi el autor puede seguir escribiendo
% «→» comodamente en el YAML.
\usepackage{newunicodechar}
\newunicodechar{→}{\ensuremath{\rightarrow}}
% Mismo problema con los simbolos de marca y vinetas: el «✓» de «una palomita ✓»
% salia como cuadrito vacio en el PDF de todas las guias que lo usaban.
\usepackage{pifont}
\newunicodechar{✓}{\ding{51}}
\newunicodechar{✗}{\ding{55}}
\newunicodechar{✔}{\ding{52}}
\newunicodechar{•}{\textbullet}
\newunicodechar{≥}{\ensuremath{\geq}}
\newunicodechar{≤}{\ensuremath{\leq}}

\setmainfont{Helvetica Neue}
\setsansfont{Helvetica Neue}
\setlength{\parindent}{0pt}
\setlength{\parskip}{6pt}
% Sin particion de palabras: en 6.º-7.º una palabra cortada por guion al final
% de linea («Comuni-cacion») frena la lectura mucho mas que una linea corta.
\hyphenpenalty=10000
\exhyphenpenalty=10000
\pretolerance=2000
\tolerance=3000
\setlength{\emergencystretch}{3em}
\linespread{1.10}
\renewcommand{\arraystretch}{1.25}
\setlist{leftmargin=5mm,itemsep=1mm,topsep=1mm}
\newcolumntype{Y}{>{\RaggedRight\arraybackslash}X}

\definecolor{milcMagenta}{HTML}{E6007E}
\definecolor{milcVino}{HTML}{5A0038}
\definecolor{milcTurquesa}{HTML}{0093A5}
\definecolor{milcVerde}{HTML}{58A923}
\definecolor{milcMostaza}{HTML}{E5B400}
\definecolor{milcOcre}{HTML}{B86600}
\definecolor{milcNegro}{HTML}{241816}
\definecolor{milcGris}{HTML}{F7F7F4}
\definecolor{milcLinea}{HTML}{E8E2DD}
\definecolor{milcGradePrimary}{HTML}{7C3AED}
\definecolor{milcGradeDark}{HTML}{4C1D95}
\definecolor{milcGradeSoft}{HTML}{EDE9FE}
\definecolor{milcGradeAccent}{HTML}{CFE7FF}
\definecolor{milcGradeText}{HTML}{FFFFFF}
\color{milcNegro}

\pagestyle{fancy}
\fancyhf{}
% Cabecera de dos lineas: identidad de la guia arriba y, a la derecha y en
% gris, la estacion en curso (\markright desde \milcestacion). Antes de la
% primera estacion la segunda linea va vacia; el \strut mantiene la altura
% para que la linea de arriba no baile entre paginas.
\lhead{\small Tecnología e Informática · Grado 9\\[-1pt]{\footnotesize\strut}}
\rhead{\small Guía 17 · MILC v3\\[-1pt]{\footnotesize\strut\color{milcNegro!65}\rightmark}}
\cfoot{\small\thepage}
\renewcommand{\headrulewidth}{0pt}

% Sobre mostaza (#E5B400) y verde (#58A923) el texto va en milcNegro; sobre el
% resto de la paleta, en blanco. Se usa dentro de fonttitle/fontupper, donde un
% \color posterior manda sobre coltitle.
\newcommand{\milctintasobre}[1]{%
  \ifboolexpr{test{\ifstrequal{#1}{milcMostaza}} or test{\ifstrequal{#1}{milcVerde}}}%
    {\color{milcNegro}}{\color{white}}}

\newtcolorbox{titlebox}[1]{
  enhanced,colback=#1,colframe=#1,arc=7mm,boxrule=0pt,
  left=7mm,right=7mm,top=3mm,bottom=3mm,
  before skip=8pt,after skip=8pt,
  fontupper=\sffamily\bfseries\Large\color{white}
}

\newtcolorbox{softbox}[3]{
  enhanced,breakable,pad at break=3mm,
  colback=#2,colframe=#1,borderline west={4pt}{0pt}{#1},
  arc=2mm,boxrule=.4pt,left=4mm,right=4mm,top=3mm,bottom=3mm,
  before skip=6pt,after skip=8pt,title={#3},coltitle=white,
  colbacktitle=#1,fonttitle=\sffamily\bfseries\milctintasobre{#1},fontupper=\RaggedRight,
  attach boxed title to top left={xshift=3mm,yshift=-2mm},
  boxed title style={arc=2mm, boxrule=0pt}
}

\newtcolorbox{drawbox}[2]{
  enhanced,breakable,pad at break=4mm,
  colback=white,colframe=#1,arc=4mm,boxrule=1pt,
  left=5mm,right=5mm,top=4mm,bottom=4mm,before skip=8pt,after skip=8pt,
  title={#2},coltitle=white,colbacktitle=#1,fonttitle=\sffamily\bfseries\milctintasobre{#1},
  fontupper=\RaggedRight,
  attach boxed title to top left={xshift=4mm,yshift=-2mm},
  boxed title style={arc=3mm, boxrule=0pt}
}

\newtcolorbox{infoband}[1]{
  enhanced,breakable,pad at break=2mm,colback=#1!8,colframe=#1!50,arc=2mm,boxrule=.5pt,
  left=4mm,right=4mm,top=2mm,bottom=2mm,before skip=4pt,after skip=4pt,
  fontupper=\small\RaggedRight
}

% Puente narrativo entre secciones (contrato editorial Conéctate).
% Da continuidad: "Acabas de X. Ahora vas a Y porque Z."
% Visualmente: banda izquierda en color del grado, texto en itálica pequeña.
\newtcolorbox{puentebox}{
  enhanced,breakable,pad at break=2mm,colback=milcGradeSoft,colframe=milcGradeDark,
  borderline west={3pt}{0pt}{milcGradeDark},
  arc=0mm,boxrule=0pt,left=5mm,right=4mm,top=2.5mm,bottom=2.5mm,
  before skip=4pt,after skip=8pt,
  fontupper=\itshape\small\color{milcNegro}\RaggedRight
}
% La flecha va en modo matematico a proposito: Helvetica Neue NO tiene el glifo
% U+2192, asi que \textrightarrow se compilaba a NADA («Missing character: There
% is no → in font Helvetica Neue») y el puente salia sin flecha. En modo
% matematico la flecha viene de la fuente matematica y si se dibuja.
\newcommand{\puente}[1]{%
  \begin{puentebox}%
    \textcolor{milcGradeDark}{$\rightarrow$}\hspace{2mm}#1%
  \end{puentebox}%
}

\newcommand{\linead}[1]{\noindent\color{milcLinea}\rule{#1}{0.5pt}\color{milcNegro}}
\newcommand{\respuesta}[1]{\par\vspace{2mm}\foreach \n in {1,...,#1}{\noindent\color{milcLinea}\rule{\linewidth}{0.5pt}\par\vspace{5mm}}\color{milcNegro}}
\newcommand{\checkbox}{\textcolor{milcGradeDark}{\fbox{\phantom{x}}}\hspace{5pt}}

% ─── Listas aireadas para las guías socioemocionales ────────────────────
% Componer las verificaciones como «\checkbox texto\\» deja el texto que salta
% de línea debajo del cuadrito y apelmaza el bloque. Estos entornos alinean la
% continuación y airean los ítems. Los usa Territorio Interior; cualquier
% familia puede hacerlo.
\newlist{checklista}{itemize}{1}
\setlist[checklista]{
  label=\textcolor{milcGradeDark}{\fbox{\phantom{x}}},
  leftmargin=9mm, labelsep=3.2mm, itemsep=2.4mm,
  topsep=2.5mm, parsep=0pt, partopsep=0pt, before=\RaggedRight
}
\newlist{pasoslista}{enumerate}{1}
\setlist[pasoslista]{
  label=\textcolor{milcGradeDark}{\sffamily\bfseries\arabic*.},
  leftmargin=8mm, labelsep=3mm, itemsep=2.8mm,
  topsep=1mm, parsep=0pt, partopsep=0pt, before=\RaggedRight
}

\newcommand{\phasecard}[4]{
\begin{tcolorbox}[enhanced,colback=#3,colframe=#2,arc=3mm,boxrule=.5pt,height=3.05cm,valign=center,halign=center,left=1.5mm,right=1.5mm,top=2mm,bottom=2mm]
{\sffamily\bfseries\fontsize{22}{24}\selectfont\textcolor{#2}{#1}}\par
{\sffamily\bfseries\small #4}
\end{tcolorbox}
}

% ── Piezas del rediseno 2026 (legibilidad 6.º-11.º) ───────────────────
% Banda de estacion: reemplaza el titlebox macizo. Titulo a la izquierda,
% dato practico (tiempo/modalidad) a la derecha, en una sola linea.
% Nunca huerfana: pide 9 lineas libres (o salta de pagina rellenando con
% \vfill) y prohibe el salto justo despues de la banda. Nueve y no seis:
% la caja partible que sigue necesita ~80pt (titulo adosado, relleno y dos
% lineas) para arrancar en la misma pagina; con menos, tcolorbox la manda
% entera a la siguiente y la banda se queda sola. El penalty 10000 va
% en `after`, ANTES del espacio posterior: un `after skip` (aunque sea 0pt)
% es un \vskip tras la caja, y ese glue es un punto de corte legal que un
% \nopagebreak escrito despues de \end{tcolorbox} ya no cierra. Cada banda
% deja marcador PDF y actualiza la cabecera derecha.
\newcounter{milcestacion}
\newcommand{\milcestacion}[2]{%
\Needspace*{9\baselineskip}%
\stepcounter{milcestacion}%
\pdfbookmark[1]{#2}{milcestacion\themilcestacion}%
\markright{#2}%
\begin{tcolorbox}[enhanced,colback=#1,colframe=#1,arc=2mm,boxrule=0pt,
  left=5mm,right=5mm,top=2.6mm,bottom=2.6mm,before skip=11pt,
  after={\par\nopagebreak\vskip7pt}]
{\sffamily\bfseries\fontsize{15}{18}\selectfont\milctintasobre{#1}#2}
\end{tcolorbox}}

% Tarjeta de paso numerada: sustituye las tablas «Pilar / Aplicacion», que
% eran formato de consulta docente, no de estudiante. El numero va en un
% circulo sobre el borde para que la secuencia se lea de un vistazo.
\newcommand{\milcpaso}[3]{%
\Needspace{4\baselineskip}%
\begin{tcolorbox}[enhanced,colback=white,colframe=milcLinea,arc=1.5mm,boxrule=.9pt,
  left=14mm,right=4mm,top=3mm,bottom=3mm,before skip=4pt,after skip=4pt,
  fontupper=\RaggedRight,
  overlay={\node[anchor=north west,circle,fill=#2,text=white,inner sep=0pt,
    minimum size=7.4mm,font=\sffamily\bfseries\small]
    at ([xshift=3.6mm,yshift=-3.2mm]frame.north west) {#1};}]
#3
\end{tcolorbox}}

% Casilla marcable de tamano util con lapiz (la anterior era \fbox{\phantom{x}},
% de alto variable segun el texto que la seguia).
\newcommand{\milcmarca}{\framebox[3.8mm]{\rule{0pt}{3.4mm}}}

% Fila marcable de lista suelta (checks de la estacion 1).
\newcommand{\milccheckrow}[1]{%
\par\vspace{1.8mm}\noindent
\makebox[6.4mm][l]{\raisebox{-.55mm}{\textcolor{milcNegro}{\milcmarca}}}%
\parbox[t]{\dimexpr\linewidth-6.4mm\relax}{\RaggedRight #1}\par}

% Celda marcable dentro de la rubrica (tabla de autoevaluacion). Misma
% sangria francesa, pero pensada para vivir dentro de una columna X.
\newcommand{\milcopcion}[1]{%
\makebox[5.2mm][l]{\raisebox{-.55mm}{\textcolor{milcNegro}{\milcmarca}}}%
\parbox[t]{\dimexpr\linewidth-5.2mm\relax}{\RaggedRight #1}}

% ── Recursos visuales de la guía ──────────────────────────────────────
% Declarados en el bloque `recursos:` del YAML y emitidos por el builder.
% \guiaFigura   → imagen rasterizada o PDF (foto, carta celeste, montaje).
% \guiaDiagrama → fuente TikZ/circuitikz (.tex) incrustada como vector:
%                 es la vía para circuitos y esquemáticos editables.
% [#1] = texto alternativo (alt) · #2 = ruta relativa al .tex · #3 = pie
% El alt llega al PDF etiquetado (/Alt) y lo lee el lector de pantalla.
\newcommand{\guiaFigura}[3][]{%
\begin{center}
\includegraphics[width=0.88\linewidth,height=0.38\textheight,keepaspectratio,alt={#1}]{#2}\\[1.5mm]
{\footnotesize\itshape\textcolor{milcNegro!75}{#3}}
\end{center}
}

\newcommand{\guiaDiagrama}[2]{%
\begin{center}
\input{#1}\\[1.5mm]
{\footnotesize\itshape\textcolor{milcNegro!75}{#2}}
\end{center}
}

\begin{document}
\thispagestyle{empty}

% ── Portada ───────────────────────────────────────────────────────────
% Antes: un rectangulo de color a pagina completa. Se veia bien en pantalla,
% pero en la fotocopiadora del colegio se comia el toner y salia gris sucio.
% Ahora la banda ocupa el tercio superior y el resto va en blanco.
\begin{tikzpicture}[remember picture,overlay]
  \path[fill=milcGradePrimary]
    (current page.north west) rectangle ([yshift=-10.6cm]current page.north east);

  \node[anchor=north west,text=milcGradeText] at ([xshift=2.0cm,yshift=-1.55cm]current page.north west)
    {\sffamily\bfseries\fontsize{12}{14}\selectfont GRADO};
  \node[anchor=north west,text=milcGradeText,opacity=.85] at ([xshift=2.0cm,yshift=-2.35cm]current page.north west)
    {\sffamily\bfseries\fontsize{8.5}{10}\selectfont SERIE GUÍAS · TECNOLOGÍA E INFORMÁTICA};
  \node[anchor=north east,text=milcGradeText,opacity=.85,align=right] at ([xshift=-2.0cm,yshift=-1.55cm]current page.north east)
    {\sffamily\bfseries\fontsize{9}{12}\selectfont I.E. Sor María Juliana\\Cartago, Valle del Cauca};

  \node[anchor=west,text=milcGradeText] (gradonum) at ([xshift=1.85cm,yshift=-7.1cm]current page.north west)
    {\sffamily\bfseries\fontsize{112}{112}\selectfont 9};
  % El circulito de grado (°) solo aplica a las guias de grado («11°»); en
  % Bebras y Semillero el numero grande es un momento/modulo, no un grado.
  % Se posiciona relativo al nodo `gradonum`, no en coordenadas absolutas.
  \draw[milcGradeText,line width=1.6pt]
    ([xshift=2.0mm,yshift=-4.5mm]gradonum.north east) circle (.15cm);

  \node[anchor=west,text=milcGradeText,text width=11.4cm,align=left] at ([xshift=1.5cm,yshift=0.5cm]gradonum.east)
    {
     {\sffamily\bfseries\fontsize{9}{11}\selectfont\MakeUppercase{Período 2 · Diseño editorial digital}}\\[3mm]
     {\sffamily\bfseries\fontsize{21}{25}\selectfont\setlength{\baselineskip}{27pt}Accesibilidad lectora ---\\[4pt]quién no puede\\[4pt]leer mi revista}\\[3.5mm]
     {\sffamily\bfseries\fontsize{15}{18}\selectfont Guía 17-9-TIC}
    };
\end{tikzpicture}

\vspace*{9.4cm}
\vfill

{\sffamily\bfseries\fontsize{8}{10}\selectfont\textcolor{milcGradeDark}{LO QUE TE LLEVAS HOY}}

\begin{tcolorbox}[enhanced,colback=white,colframe=milcNegro,arc=2mm,boxrule=1.6pt,
  left=5mm,right=5mm,top=4mm,bottom=4mm,before skip=3pt,after skip=12pt,
  fontupper=\RaggedRight\fontsize{11}{15.5}\selectfont]
Un spread de tu revista rediseñado para que sea accesible: contraste de texto y fondo verificado, cuerpo de letra suficiente, texto alternativo en cada imagen informativa y una jerarquía de títulos que un lector de pantalla pueda recorrer. Va con la bitácora de antes y después, donde cada cambio dice qué barrera quita y a quién deja entrar.
\end{tcolorbox}

\begin{tcolorbox}[enhanced,colback=milcGris,colframe=milcGris,arc=2mm,boxrule=0pt,
  left=5mm,right=5mm,top=3.5mm,bottom=3.5mm,after skip=12pt]
{\sffamily\bfseries\fontsize{8}{10}\selectfont\textcolor{milcNegro!55}{ESTA GUÍA ES DE}}\\[3mm]
\begin{tabularx}{\linewidth}{X p{3.2cm} p{3.2cm}}
\linead{\linewidth} & \linead{3.0cm} & \linead{3.0cm}\\[-1mm]
{\fontsize{8.5}{10}\selectfont\textcolor{milcNegro!55}{Nombre completo}} &
{\fontsize{8.5}{10}\selectfont\textcolor{milcNegro!55}{Grupo}} &
{\fontsize{8.5}{10}\selectfont\textcolor{milcNegro!55}{Fecha}}\\
\end{tabularx}
\end{tcolorbox}

\vfill

\noindent\textcolor{milcLinea}{\rule{\linewidth}{1pt}}\\[2mm]
{\sffamily\bfseries\fontsize{9.5}{12}\selectfont Autor: Dr. Álvaro Cárdenas Orozco}\\[1mm]
{\sffamily\fontsize{8.5}{11}\selectfont\textcolor{milcNegro!55}{Metodología MILC · apertura ancestral · pensamiento computacional · triángulo Dussel–estoicismo–Floridi}}

\newpage

\section*{Apertura: Accesibilidad lectora --- quién NO puede leer mi revista y por qué}

\begin{softbox}{milcGradeDark}{milcGradeSoft}{Saber ancestral}
Colombia tardó medio siglo en decir bien una sola cosa. La \textbf{Ley 324 de 1996} reconoció la lengua de los sordos colombianos. Pero la llamó \textbf{«lengua manual»}, como si fuera un truco con las manos. Hubo que esperar a la \textbf{Ley 982 de 2005} para que la ley dijera lo correcto: la lengua de señas es la \textbf{lengua natural} de la comunidad sorda. La diferencia entre las dos leyes no es de palabras. Es la diferencia entre «este niño tiene un problema» y «este niño tiene un idioma que la escuela no habla». Fíjate en lo que eso cambia. En la primera versión, quien debe adaptarse es la persona. En la segunda, quien debe adaptarse es la institución. En un salón del Valle del Cauca esa diferencia decide algo concreto: si un estudiante sordo entra a la clase o se queda mirándola. La \textbf{cara de exclusión}: la exclusión no necesita mala intención. Basta con no hablar el idioma del otro y no haberse dado cuenta.\par\smallskip{\footnotesize\textcolor{milcNegro!70}{\textbf{Fuente.} Congreso de la República de Colombia. (2005, 2 de agosto). \emph{Ley 982 de 2005. Por la cual se establecen normas tendientes a la equiparación de oportunidades para las personas sordas y sordociegas}. Diario Oficial.}}
\end{softbox}

\begin{softbox}{milcTurquesa}{milcGris}{Saber contemporáneo}
La accesibilidad de una pieza se mide con cuatro ejes. \textbf{Perceptible}: que se pueda ver u oír, con contraste suficiente y con texto alternativo en las imágenes. \textbf{Operable}: que se pueda recorrer sin ratón y sin depender de un gesto fino. \textbf{Comprensible}: lenguaje claro y comportamiento previsible. \textbf{Robusto}: que funcione con las herramientas que usa la gente, incluidos los lectores de pantalla. Hay un nivel de referencia que sirve como piso realista. Un contraste de al menos 4,5 a 1 entre texto y fondo. Un cuerpo de al menos 16 píxeles en digital. Un texto alternativo por imagen informativa. Un solo título principal por página. Y una pregunta que ordena todo lo anterior: quién queda afuera con cada decisión que tomo. Sin esa pregunta, la accesibilidad se queda en buena intención.
\end{softbox}

\begin{softbox}{milcMostaza}{milcGris}{Pregunta-puente}
\textit{La ley pasó de decir «lengua manual» a decir «lengua natural», y con eso cambió quién tenía que adaptarse. En tu revista, ¿quién se está adaptando a quién?}
\end{softbox}

\begin{softbox}{milcVerde}{milcGris}{Saber y saber hacer}
Al terminar podrás: (1) \textbf{identificar} qué barreras encuentra tu revista simulando cuatro condiciones de lectura distintas; (2) \textbf{explicar} los cuatro ejes de accesibilidad y el piso mínimo de cada uno; (3) \textbf{aplicar} esos criterios a un spread real y documentar qué barrera quita cada cambio.
\end{softbox}



\small
\begin{tabularx}{\textwidth}{>{\bfseries}p{3.0cm}Y}
\rowcolor{milcGradePrimary!12}Autor & Dr. Álvaro Cárdenas Orozco\\
DBA & Producción de piezas editoriales digitales con criterio estético y ético (MEN, T\&I)\\
Referentes & Tipografía colombiana · Diseño centrado en accesibilidad · Floridi\\
Producto final & Un spread de tu revista rediseñado para que sea accesible: contraste de texto y fondo verificado, cuerpo de letra suficiente, texto alternativo en cada imagen informativa y una jerarquía de títulos que un lector de pantalla pueda recorrer. Va con la bitácora de antes y después, donde cada cambio dice qué barrera quita y a quién deja entrar.\\
\end{tabularx}
\normalsize

\puente{En las sesiones anteriores armaste páginas, elegiste tipografía y color y montaste spreads. Hoy compruebas quién no puede leerlos. En la sesión 8 corriges, y esta revisión te va a dar la mitad de la lista.}

\milcestacion{milcGradeDark}{Ruta MILC: así vamos a aprender}
\begin{tabularx}{\textwidth}{>{\centering\arraybackslash}X>{\centering\arraybackslash}X>{\centering\arraybackslash}X>{\centering\arraybackslash}X}
\phasecard{1}{milcVerde}{milcGris}{Escucha\\[-1mm]\small Observo, escucho y pregunto.} &
\phasecard{2}{milcTurquesa}{milcGris}{Entender\\[-1mm]\small Sistematizo y ordeno.} &
\phasecard{3}{milcMagenta}{milcGris}{Hacer\\[-1mm]\small Construyo y aplico.} &
\phasecard{4}{milcVino}{milcGris}{Cierre\\[-1mm]\small Evalúo y reflexiono.}
\end{tabularx}


\puente{Antes de la teoría vas a hacer la prueba en tu propio cuerpo. Cuatro simulaciones de dos minutos cada una, sobre tu revista, sin tocar nada todavía.}

\milcestacion{milcVerde}{Estación 1 de 3 · Escucha}

\begin{softbox}{milcVerde}{milcGris}{Escena para escuchar}
\textbf{Actividad 1 · IDENTIFICA --- Cuatro barreras en tu revista} (15 min · individual). (1) Baja visión: aleja la revista o entrecierra los ojos y anota qué texto dejas de leer. (2) Daltonismo: pásale un simulador o imagina la pieza en escala de grises y anota qué información se pierde. (3) Sin imágenes: tapa las fotos y comprueba si el texto solo sigue explicando. (4) Solo audio: pídele a alguien que te lea la página en voz alta y anota qué no se entiende. (5) Marca cuál barrera dejó afuera a más gente.
\end{softbox}

{\sffamily\bfseries\fontsize{8}{10}\selectfont\textcolor{milcNegro!55}{MARCA CUANDO LO HAYAS HECHO}}
\milccheckrow{Hice las cuatro simulaciones sobre mi propia revista.}
\milccheckrow{Anoté qué información se pierde en cada una.}
\milccheckrow{Marqué la barrera que deja afuera a más gente.}

\begin{infoband}{milcVerde}


\textbf{Lo que va al cuaderno (Actividad 1 · IDENTIFICA).} \textbf{Título:} «Actividad 1 --- Cuatro barreras en mi revista». \textbf{Formato:} tabla de 4 filas y 3 columnas (barrera / qué se pierde / a quién deja afuera), con marca en la peor. \textbf{Extensión:} un tercio de página. \textbf{Sabes que terminaste cuando} cada fila nombra qué información concreta se pierde.
\end{infoband}

\puente{Ya sabes qué falla. Ahora vas a ordenarlo en los cuatro ejes y a fijar el piso mínimo de cada uno, para no discutir por gusto.}

\milcestacion{milcTurquesa}{Estación 2 de 3 · Entender}

\begin{softbox}{milcTurquesa}{milcGris}{Lo que vas a entender}
\textbf{Actividad 2 · EXPLICA --- Los cuatro ejes y su piso} (20 min · parejas). (1) Con tu pareja, escriban los cuatro ejes con una frase propia cada uno. (2) Ubiquen en ellos las cuatro barreras que encontraron. (3) Escriban el piso mínimo de cada eje, con su número cuando lo tenga. (4) Elijan un cambio que arregle dos ejes a la vez y expliquen por qué funciona doble.
\end{softbox}

\milcpaso{1}{milcTurquesa}{\textbf{Los cuatro ejes.} Perceptible: que se pueda ver u oír. Operable: que se pueda recorrer sin ratón. Comprensible: lenguaje claro y comportamiento previsible. Robusto: que funcione con lectores de pantalla.}
\milcpaso{2}{milcTurquesa}{\textbf{El piso mínimo tiene números.} Contraste de al menos 4,5 a 1 entre texto y fondo. Cuerpo de al menos 16 píxeles en digital. Un texto alternativo por imagen informativa. Un solo título principal por página.}
\milcpaso{3}{milcTurquesa}{\textbf{El texto alternativo describe la información, no la imagen.} «Foto» no sirve. «Gráfico: el consumo baja a la mitad después de marzo» sí, porque quien no ve la imagen recibe lo mismo que quien la ve.}
\milcpaso{4}{milcTurquesa}{\textbf{Lo que ayuda a uno suele ayudar a todos.} Un contraste alto sirve a quien tiene baja visión y a quien lee al sol. Los subtítulos sirven a quien no oye y a quien viaja en bus. Casi nunca es un arreglo para pocos.}

\begin{drawbox}{milcTurquesa}{Los cuatro ejes y su piso}
\textbf{Perceptible:} contraste $\geq$ 4,5 a 1; texto alternativo. \\[1mm] \textbf{Operable:} recorrible sin ratón. \\[1mm] \textbf{Comprensible:} lenguaje claro, comportamiento previsible. \\[1mm] \textbf{Robusto:} un solo título principal, estructura legible por lector de pantalla.
\end{drawbox}

\begin{softbox}{milcMostaza}{milcGris}{Errores comunes y su corrección}
(1) \textbf{Texto gris claro sobre blanco}: se ve elegante en pantalla y desaparece al sol. (2) \textbf{Información solo por color}: si el rojo significa «mal», quien no distingue rojo no lee nada. (3) \textbf{Texto alternativo que dice «imagen»}: no aporta, y ocupa el lugar del que sí aportaría. (4) \textbf{Varios títulos principales}: el lector de pantalla pierde la estructura. (5) \textbf{Dejar la accesibilidad para el final}: rehacer cuesta más que decidirlo antes.
\end{softbox}

\begin{infoband}{milcTurquesa}


\textbf{Lo que va al cuaderno (Actividad 2 · EXPLICA).} \textbf{Título:} «Actividad 2 --- Los cuatro ejes y su piso». \textbf{Formato:} los cuatro ejes con frase propia, las cuatro barreras ubicadas y el piso mínimo de cada eje con su número. \textbf{Extensión:} media página. \textbf{Sabes que terminaste cuando} tienes escrito un cambio que arregla dos ejes a la vez.
\end{infoband}

\puente{Tienes el criterio. Toca rediseñar un spread de verdad y anotar qué barrera quita cada cambio.}

\milcestacion{milcMagenta}{Estación 3 de 3 · Hacer}

\begin{softbox}{milcMagenta}{milcGris}{Cómo vas a construir tu producto}
\textbf{Actividad 3 · APLICA --- Rediseña un spread} (30 min · individual). (1) Elige un spread de tu revista, el más débil de la prueba anterior. (2) Verifica el contraste de texto y fondo y súbelo si no llega al piso. (3) Sube el cuerpo de letra hasta el mínimo. (4) Escribe el texto alternativo de cada imagen informativa, en una o dos frases con la información. (5) Deja un solo título principal y ordena los subtítulos debajo. (6) Anota cada cambio en la bitácora, diciendo qué barrera quita y a quién deja entrar.
\end{softbox}

\milcpaso{1}{milcMagenta}{\textbf{Tu entrega tiene dos piezas.} El spread rediseñado, con el antes y el después visibles. Y la bitácora, con un renglón por cambio.}
\milcpaso{2}{milcMagenta}{\textbf{Cada renglón de la bitácora nombra a alguien.} «Subí el contraste» describe. «Subí el contraste para que se lea al sol y con baja visión» explica. La segunda forma es la que sirve cuando alguien pregunte por qué.}
\milcpaso{3}{milcMagenta}{\textbf{El texto alternativo se escribe pensando en voz alta.} Léelo como si se lo estuvieras contando por teléfono a alguien que no ve la página. Si esa persona se queda igual de informada, está bien escrito.}
\milcpaso{4}{milcMagenta}{\textbf{Un cambio suele arreglar dos cosas.} Subir el cuerpo de letra mejora lo perceptible y de paso lo comprensible, porque obliga a cortar texto. Anótalo cuando pase: es la parte que se aprende.}

\begin{drawbox}{milcMagenta}{Antes de entregar}
\checkbox Contraste de texto y fondo verificado y por encima del piso. \\[1mm] \checkbox Cuerpo de letra en el mínimo o por encima. \\[1mm] \checkbox Texto alternativo en cada imagen informativa. \\[1mm] \checkbox Un solo título principal y subtítulos ordenados. \\[1mm] \checkbox Bitácora con antes y después por cada cambio. \\[1mm] \checkbox Cada renglón dice a quién deja entrar.
\end{drawbox}

\begin{softbox}{milcMagenta}{milcGris}{Plantilla de trabajo}
\textbf{Cambio.} \textit{Qué moví.} \\[1mm] \textbf{Eje.} \textit{Perceptible, operable, comprensible o robusto.} \\[1mm] \textbf{Antes.} \textit{Cómo estaba.} \\[1mm] \textbf{Después.} \textit{Cómo quedó.} \\[1mm] \textbf{A quién deja entrar.} \textit{Dicho con nombre de la situación.}
\end{softbox}

\begin{infoband}{milcMagenta}


\textbf{Lo que va al cuaderno (Actividad 3 · APLICA).} \textbf{Título:} «Actividad 3 --- Rediseña un spread». \textbf{Formato:} el boceto del antes y el después, y la bitácora con un renglón por cambio, cada uno con su eje y a quién deja entrar. \textbf{Extensión:} media página. \textbf{Sabes que terminaste cuando} ningún renglón de la bitácora se queda en describir el cambio.
\end{infoband}

\puente{Lo que entregas hoy: un spread rediseñado y su bitácora. La prueba de calidad: cada cambio nombra a quién deja entrar, no solo qué se movió.}

\milcestacion{milcMostaza}{Producto de la sesión}
\textbf{Un spread accesible + bitácora donde cada cambio dice a quién deja entrar}.
\begin{softbox}{milcMostaza}{milcGris}{Criterios de calidad}
(1) El contraste de texto y fondo está verificado y por encima del piso. (2) El cuerpo de letra alcanza el mínimo de la guía. (3) Cada imagen informativa tiene texto alternativo con la información, no la palabra «imagen». (4) Hay un solo título principal y los subtítulos están ordenados debajo. (5) Cada renglón de la bitácora dice qué barrera quita y a quién deja entrar.
\end{softbox}

\puente{Antes de cerrar, mira lo que hiciste desde cinco lados. Nadie diseña una barrera a propósito, y sin embargo casi todas las páginas tienen alguna.}

\milcestacion{milcVino}{Evaluación liberadora · cinco dimensiones}

\begin{softbox}{milcVino}{milcGris}{Sentido de la evaluación}
Evaluar no es solo revisar una nota. Es reconocer qué aprendí, qué cambió en mí, qué emociones manejé, cómo me relaciono con la ciudad y la comunidad, y qué le dejaré a quien viene detrás.
\end{softbox}

\textbf{Mi reflexión liberadora en cinco dimensiones.} Completa cada frase iniciada:

\small
\begin{tabularx}{\textwidth}{>{\bfseries\RaggedRight\arraybackslash}p{3.4cm}Y>{\RaggedRight\arraybackslash}p{4.6cm}}
\rowcolor{milcVino}\color{white}Dimensión & \color{white}\bfseries Mi respuesta (frame starter) & \color{white}\bfseries Algo que descubrí\\
1. Personal & Lo que más cambió en mí fue \linead{6.5cm} & \linead{4.2cm}\\
2. Emocional & La emoción más fuerte fue \linead{2.5cm} y la regulé \linead{3cm} & \linead{4.2cm}\\
3. Ciudadana & En democracia esto importa porque \linead{6cm} & \linead{4.2cm}\\
4. Local (Cartago) & En mi barrio sería útil para \linead{6.5cm} & \linead{4.2cm}\\
5. Intergeneracional & A mi abuela/abuelo le diría \linead{6.5cm} & \linead{4.2cm}\\
\end{tabularx}
\normalsize

\begin{infoband}{milcVino}
\textbf{Para la próxima sustentación.} Vuelve a esta tabla en seis meses. Verás que las respuestas cambian — no porque lo aprendido cambie, sino porque tú cambias. La evaluación liberadora es una conversación contigo a través del tiempo.
\end{infoband}


\puente{Para terminar, tres citas verificadas sobre escuchar y no excluir: Dussel sobre la conciencia ética como capacidad de oír al otro, Marco Aurelio sobre el enjambre y la abeja, Floridi sobre no reducir a nadie a instrumento.}

\milcestacion{milcGradeDark}{Triángulo de pensamiento · cierre filosófico}

\begin{softbox}{milcVino}{milcGris}{Dussel · lente del nosotros}
\textit{«Llamamos conciencia ética la capacidad que se tiene de escuchar la voz del otro, palabra transontológica que irrumpe desde más allá del sistema vigente. Puede que la protesta justa del otro ponga en cuestión los principios morales del sistema.»}

{\footnotesize\itshape\color{milcNegro!70}--- Enrique Dussel · Filosofía de la liberación (1977), §2.6.2.2}

La ley de 1996 no oyó bien y llamó «manual» a una lengua. La de 2005 corrigió. Entre las dos hay nueve años y una comunidad insistiendo, que es exactamente lo que significa que una protesta ponga en cuestión lo vigente.

\textbf{¿Qué me han dicho sobre mi trabajo que no quise oír porque me obligaba a rehacerlo?}
\end{softbox}
\linead{\linewidth}\\[1mm]
\linead{\linewidth}

\begin{softbox}{milcMostaza}{milcGris}{Estoicismo · Marco Aurelio · Meditaciones VI, 54 (c. 175 d.C.) · cuidado interior}
\textit{«Lo que no es bueno para el enjambre, tampoco es conducente para la abeja.»}

{\footnotesize\itshape\color{milcNegro!70}--- Marco Aurelio · Meditaciones VI, 54 (c. 175 d.C.)}

Y al revés también: lo que se hace por unos pocos termina sirviéndole a todos. El contraste alto lo agradece quien tiene baja visión y quien lee en la calle al mediodía. Casi nada de esto es un arreglo para pocos.

\textbf{¿Qué arreglo hice pensando en alguien concreto y terminó mejorando la pieza entera?}
\end{softbox}
\linead{\linewidth}\\[1mm]
\linead{\linewidth}

\begin{softbox}{milcTurquesa}{milcGris}{Floridi · lente de la infoesfera}
\textit{«En una esfera pública plural, los demás no pueden reducirse a instrumentos, y hacen falta contención y respeto. (trad. propia, abreviada)»}

{\footnotesize\itshape\color{milcNegro!70}--- The Onlife Initiative (ed. Luciano Floridi) · The Onlife Manifesto (2015), § 4.3}

Publicar es entrar a un espacio donde no todos leen igual. Diseñar suponiendo un único tipo de lector no es un error técnico: es haber decidido, sin decirlo, para quién existe la pieza.

\textbf{¿A qué lector le estoy diseñando por defecto, y a quién estoy dejando fuera sin haberlo decidido?}
\end{softbox}
\linead{\linewidth}\\[1mm]
\linead{\linewidth}

\begin{infoband}{milcNegro}
\textbf{Nota para el docente.} Las tres son citas textuales y llevan su referencia debajo. Puedes remitir a la obra en clase.
\end{infoband}

\milcestacion{milcVino}{Mi autoevaluación · marca una casilla por fila}

{\small
\setlength{\extrarowheight}{2.2pt}
\begin{tabularx}{\textwidth}{>{\RaggedRight\arraybackslash\bfseries}p{3.0cm}YYY}
\rowcolor{milcVino}\color{white}\bfseries Criterio & \color{white}\bfseries Lo logré & \color{white}\bfseries En proceso & \color{white}\bfseries Necesito apoyo\\[.6mm]
Comprensión del tema & \milcopcion{Explico las ideas con mis palabras.} & \milcopcion{Reconozco las ideas, falta precisión.} & \milcopcion{Necesito repasar lo básico.}\\[.8mm]
\rowcolor{milcGris}Pensamiento computacional & \milcopcion{Usé los pilares al armar mi producto.} & \milcopcion{Usé algunos pilares.} & \milcopcion{Necesito releer la estación 2.}\\[.8mm]
Aplicación y producto & \milcopcion{Mi producto cumple los criterios.} & \milcopcion{Está armado, con ajustes pendientes.} & \milcopcion{Aún está en construcción.}\\[.8mm]
\rowcolor{milcGris}Reflexión 5 dimensiones & \milcopcion{Respondí cada una con honestidad.} & \milcopcion{Respondí algunas con detalle.} & \milcopcion{Mis respuestas son muy generales.}\\[.8mm]
Triángulo de pensamiento & \milcopcion{Conecté las 3 voces con mi trabajo.} & \milcopcion{Conecté al menos una.} & \milcopcion{No las relaciono con mi proceso.}\\[.8mm]
\end{tabularx}}

\vspace{3mm}
\textbf{Hoy entendí que:}\\[1mm]
\linead{\linewidth}\\[1mm]
\linead{\linewidth}

\textbf{Mi compromiso para la próxima sesión es:}\\[1mm]
\linead{\linewidth}\\[1mm]
\linead{\linewidth}

\begin{softbox}{milcMostaza}{milcGris}{Cita de despedida · Marco Aurelio}
\textit{``El obstáculo en el camino se vuelve el camino. Lo que estorba al camino se vuelve camino.''}\\[1mm]
Cada error de hoy, bien observado, se convierte en pieza del camino que pisarás mañana.
\end{softbox}

\small
\begin{tabularx}{\textwidth}{>{\bfseries}p{3.6cm}Y}
\rowcolor{milcGradeDark!12}Nombre completo: & \linead{10cm}\\
Fecha: & \linead{4cm} \quad Curso/grupo: \linead{4cm}\\
Mi firma: & \linead{9cm}\\
\end{tabularx}
\normalsize

\begin{infoband}{milcGradeDark}
\textbf{Para la próxima sesión.} Llega con el spread rediseñado y su bitácora. Pásale la revista a alguien que use gafas y pregúntale qué le costó leer. En la sesión 8 aplicas el ojo de corrector sobre tu propio trabajo, y esta lista es la mitad del camino.
\end{infoband}

\end{document}
